\section{Pojmi, splošno}

\pojem{Izid} je en od možnih rezultatov nekega verjetnostnega poskusa (met
kovanca, met kocke, mešanje kupa kart, ipd.).
Množico vseh izidov običajno označimo z $\Omega$.
Množica nekih izidov se imenuje \pojem{dogodek}.
Dogodkom pripisujemo verjetnosti s preslikavo $P: 2^\Omega \to [0,1]$,
kjer velja $P(\varnothing) = 0$, $P(\Omega) = 1$.
Za števno družino disjunktnih dogodkov velja
\[
  P\left( \bigcup_i A_i \right) = \sum_i P(A_i).
\]
Dodatno velja $P(A\complement) = 1 - P(A)$ in formula za vključitve in
izključitve
\[
  P\!\left(\,\bigcup_{i=1}^n A_i \right)
  = \sum_{k = 1}^n {(-1)}^{k-1} \sum_{i_1, i_2, \ldots, i_k}
    P\!\left( \bigcap_j A_{i_j} \right)
\]

\pojem{Pogojna verjetnost} dogodka $A$ glede na dogodek $B$ je
\[
  P(A \such B) = \frac{P(A \presek B)}{P(B)}.
\]
Dogodka sta neodvisna, če velja $P(A \such B) = P(A)$, oziroma
$P(A \presek B) = P(A) P(B)$.
Če je družina dogodkov $H_i$ particija množice $\Omega$, lahko uporabimo formulo
za popolno verjetnost
\[
  P(A) = \sum_i P(A \such H_i) P(H_i).
\]
Poleg tega velja Bayesova formula
\[
  P(A \such B) = \frac{P(A) P(B \such A)}{P(B)}.
\]

\pojem{Slučajna spremenljiva} $X$ je funkcija $\Omega \to \R$.
Pri tem si mislimo, da krempelj našega Velikega Škampa, Boga Vsemogočnega,
izbere en izid, ki ga vstavi kot argument v funkcijo $X$; na funkciji sami po
sebi ni ničesar naključnega.
Pišemo
\[
  P(X = k) = P(X^{-1}(\{k\})).
\]
Najpomembnejše pri slučajnih spremenljivkah so njihove porazdelitve, t.j.
vrednosti $P(X = k)$ za vse možne $k$.
Porazdelitve so naštete v Bronštejnu na strani 602 in v
razdelku~\ref{sec:porazdelitve} spodaj.
Poseben tip spremenljivk so slučajni vektorji; namesto $\Omega \to \R$ so to
funkcije $\Omega \to \R^k$.
Označimo $\underline{X} = (X_1, \ldots, X_k)$ in pišemo
\[
  P(\underline{X} = \underline{x}) = P(X_1 = x_1, \ldots, X_k = x_k).
\]

\section{Reševanje nalog}

Najpomembnejši korak pri vsaki nalogi je določitev množice možnih izidov
$\Omega$.
Običajno je to neko zaporedje ali urejena $k$-terica.
Neurejenih $k$-teric se izogibaj, razen če znaš utemeljiti, zakaj ti da pravi
odgovor.
Vsakemu izidu potem prirediš verjetnost --- običajno kar vsem enako, ne pa
nujno.

\section{Verjetnostne porazdelitve}%
\label{sec:porazdelitve}

\subsection{Diskretne porazdelitve}

\begin{itemize}
\item \pojem{Binomska porazdelitev}:
  $n$-krat vržemo kovanec z verjetnostjo grba $p$.
  Če je $X$ število grbov v $n$ metih, je
  \[
	P(X = k) = \binom{n}{k} p^k q^{n-k}
	\sim \operatorname{Bin}(n, p).
  \]
\item \pojem{Negativna binomska porazdelitev}:
  Mečemo kovanec in čakamo na $m$ grbov.
  Naj bo $X$ potrebno število metov.
  Velja
  \[
	P(X=k) = \binom{k-1}{m-1} p^m q^{k-m} \sim \operatorname{NegBin}(m, p).
  \]
\item \pojem{Geometrijska porazdelitev}:
  Mečemo kovanec in čakamo na prvi grb. Če je $X$ število potrebnih metov, je
  \[
	P(X = k) = q^{k-1}p \sim \operatorname{Geom}(p).
  \]
\item \pojem{Hipergeometrijska porazdelitev}:
  Imamo posodo z $B$ belimi in $R$ rdečimi kroglicami, naključno izberemo $n \le
  N = B+R$ kroglic iz posode.
  Če je $X$ število izbranih belih kroglic, dobimo
  \[
	P(X = k) = \frac{\binom{B}{k} \binom{N-B}{n-k}}{\binom{N}{n}}
	\sim \operatorname{HiperGeom}(n, B, N).
  \]
\item \pojem{Poissonova porazdelitev}:
  Tu je
  \[
	P(X = k) = e^{-\lambda} \frac{\lambda^k}{k!} \sim \operatorname{Po}(\lambda).
  \]
\end{itemize}

\subsection{Porazdelitve slučajnih vektorjev}

\begin{itemize}
\item \pojem{Multinomska porazdelitev}:
  Imamo $r$ škatel, v katere mečemo $n$ kroglic.
  Vsako škatlo zadenemo z verjetnostjo $p_i$, kjer je $\sum_i p_i = 1$.
  Naj bo $X_i$ število kroglic v $i$-ti škatli na koncu metov.
  Tedaj
  \[
	P(X_1 = k_1, X_2 = k_2, \ldots, X_r = k_r)
	= p_1^{k_1} \ldots p_r^{k_r} \binom{n}{k_1, k_2, \ldots, k_r}.
	\sim \operatorname{Multinom}(n, \underline{p}).
  \]
\end{itemize}

\subsection{Zvezne porazdelitve}

Slučajna spremenljivka ima zvezno porazdelitev, če obstaja taka funkcija $f_X:\R
\to [0, \infty)$, da je
\[
  P(X \in (a, b])
  = \int_a^b f_X(x) dx.
\]
Taki funkciji pravimo \pojem{gostota porazdelitve}.
Znane zvezne porazdelitve so naslednje:
\begin{itemize}
\item \pojem{Normalna porazdelitev}:
  \[
	f_X(x) = \inv{\sigma \sqrt{2 \pi}} \exp\left( - \frac{{(x - \mu)}^2}{2
		\sigma^2} \right)
  \]
\item \pojem{Eksponentna porazdelitev}:
  \[
	f_X(x) = \lambda \exp(-\lambda x)
  \]
  za $x \ge 0$.
\item \pojem{Gama porazdelitev}:
  \[
	f_X(x) = \frac{\lambda^a}{\Gamma(a)} x^{a-1} \exp(-\lambda x)
  \]
  za $x \ge 0$.
\item \pojem{Enakomerna porazdelitev}:
  \[
	f_X(x) = \inv{b-a}.
  \]
\item \pojem{Beta porazdelitev}:
  \[
	f_X(x) = \inv{B(a, b)} x^{a-1} (1-x)^{b-1}
  \]
  za $0 < x < 1$.
\end{itemize}
